\paragraph{Scope of exact sampling}
The sampling rule is exact for the intended target distribution only when the supplied probabilities are correct. History-dependent processing (presence, frequency, and repetition penalties) must therefore use each verification row's own ancestry: a target row includes the path through its parent, and a leaf continuation row includes the leaf itself. The deployed processor builds these per-path histories from the tree descriptor's parent table; the stock chain combiner, which gives row $j$ the flattened history of the first $j$ drafts, is retained for chain requests only. Two differences from native single-row decoding remain: the tree path ignores min-$p$, and its top-$k$ implementation retains exactly $k$ tokens, whereas native decoding retains boundary ties. The numerical continuation result does not prove joint output-distribution equality.

\paragraph{Captured-row and sampler audit}
A passive capture of 40 tree steps on four requests, taken with the deployed sampler at presence penalty 1.0, contains 2,480 processed target and continuation rows, including padded slots. All 2,480 carry the per-path history; 82 coincide with the flattened-chain history where the two agree. An earlier capture before this correction found all 2,480 rows flattened, with node-level total variation above 0.01 on approximately 22\% of deeper decisions and maximum 0.421; that defect, its repair, and the superseded measurements are in the supplement. Every measurement here uses the corrected sampler unless stated otherwise.

The deployed target-only walk consumes target and continuation logits without a separate draft-probability tensor. Its CPU floating-point oracle is sampled 200,000 times on each of six captured steps, for 1.2 million walks. Tests compare emitted-token counts conditional on visiting a node with the supplied probability row, using an approximate chi-square $G$-test and a Bonferroni-adjusted per-step threshold of $10^{-3}$. Nodes require at least 2,000 visits; rare expected counts are pooled. No departure is detected in 48 node tests against the supplied rows. Against autoregressive rows reconstructed with the native tie-keeping top-$k$, one of 48 node tests rejects; over all 1,120 node-steps the reconstructed total variation is at most 0.028, and the three nodes above 0.01, including the rejected one, are exact ties at the top-$k$ boundary. On these rows, no departure from the autoregressive target is found beyond the stated tie convention. All five planted fault types are detected in at least one step (Table~\ref{tab:sampling-faults}). Non-rejection supports this finite check; it does not prove universal exactness.

\begin{table}[t]
\caption{Sampler-audit fault sensitivity on six captured steps. Each fault type is detected at least once; detection depends on which nodes the walk visits.}
\label{tab:sampling-faults}
\centering\small
\begin{tabular}{@{}lr@{}}
\toprule
Injected fault & Steps detected \\
\midrule
Temperature applied twice & 5/6 \\
Target rows shifted & 6/6 \\
Fixed residual draw & 4/6 \\
Unconditional acceptance & 3/6 \\
Correlated random draws & 1/6 \\
\bottomrule
\end{tabular}
\end{table}

An earlier ten-boot, 8,000-continuation output diagnostic failed its frozen temperature-control sensitivity requirement and remains uninformative. The tree sampler also seeds one process-wide generator with zero, so boots with identical request order reuse random sequences. They cannot be treated as independent sampling replicates. The captured-row audit above is informative where that output test was not.
